% =====================================================================
%  introduction.tex  --  PHASE VII, Task 14E
%
%  Section 1.  Context and problem, the positioning statement, the
%  contribution list, and the organisation paragraph.
%
%  Supersedes the corresponding blocks of novelty_and_contributions.tex
%  (Phase I), which recorded the design rationale for the reframing and
%  is retained for that purpose only.  Changes from the Phase I text:
%
%   * component (ii) is a duration forecaster benchmarked against an
%     untrained reference, not "a Seq2Seq LSTM" -- Phase V;
%   * the accept/reject loop is promoted to a named contribution
%     -- Phase V Task 12B;
%   * the decision contribution claims the paired contrast against the
%     static rule, not a saving against the status quo -- Phase V 11B;
%   * the sustainability/deployment forward references are removed;
%     those analyses are not in this paper.
%
%  Drop-in usage: % =====================================================================
%  introduction.tex  --  PHASE VII, Task 14E
%
%  Section 1.  Context and problem, the positioning statement, the
%  contribution list, and the organisation paragraph.
%
%  Supersedes the corresponding blocks of novelty_and_contributions.tex
%  (Phase I), which recorded the design rationale for the reframing and
%  is retained for that purpose only.  Changes from the Phase I text:
%
%   * component (ii) is a duration forecaster benchmarked against an
%     untrained reference, not "a Seq2Seq LSTM" -- Phase V;
%   * the accept/reject loop is promoted to a named contribution
%     -- Phase V Task 12B;
%   * the decision contribution claims the paired contrast against the
%     static rule, not a saving against the status quo -- Phase V 11B;
%   * the sustainability/deployment forward references are removed;
%     those analyses are not in this paper.
%
%  Drop-in usage: % =====================================================================
%  introduction.tex  --  PHASE VII, Task 14E
%
%  Section 1.  Context and problem, the positioning statement, the
%  contribution list, and the organisation paragraph.
%
%  Supersedes the corresponding blocks of novelty_and_contributions.tex
%  (Phase I), which recorded the design rationale for the reframing and
%  is retained for that purpose only.  Changes from the Phase I text:
%
%   * component (ii) is a duration forecaster benchmarked against an
%     untrained reference, not "a Seq2Seq LSTM" -- Phase V;
%   * the accept/reject loop is promoted to a named contribution
%     -- Phase V Task 12B;
%   * the decision contribution claims the paired contrast against the
%     static rule, not a saving against the status quo -- Phase V 11B;
%   * the sustainability/deployment forward references are removed;
%     those analyses are not in this paper.
%
%  Drop-in usage: % =====================================================================
%  introduction.tex  --  PHASE VII, Task 14E
%
%  Section 1.  Context and problem, the positioning statement, the
%  contribution list, and the organisation paragraph.
%
%  Supersedes the corresponding blocks of novelty_and_contributions.tex
%  (Phase I), which recorded the design rationale for the reframing and
%  is retained for that purpose only.  Changes from the Phase I text:
%
%   * component (ii) is a duration forecaster benchmarked against an
%     untrained reference, not "a Seq2Seq LSTM" -- Phase V;
%   * the accept/reject loop is promoted to a named contribution
%     -- Phase V Task 12B;
%   * the decision contribution claims the paired contrast against the
%     static rule, not a saving against the status quo -- Phase V 11B;
%   * the sustainability/deployment forward references are removed;
%     those analyses are not in this paper.
%
%  Drop-in usage: \input{introduction}
%  Bibliography: paper/references.bib
% =====================================================================

\section{Introduction}
\label{sec:intro}

Electric motor systems account for the largest single share of
industrial electricity use, and a substantial part of that share is
drawn while the machine is energised but producing nothing. An induction
motor left running between production runs continues to draw magnetising
current: real power falls to a fraction of its loaded value, but it does
not fall to zero, and the machine's power factor collapses. ISO~14955-1
\cite{iso14955} names this operating mode --- ready for production,
performing no value-adding function --- and studies of discrete-part
manufacturing find that such load-independent consumption frequently
exceeds the energy actually delivered to the workpiece
\cite{gutowski2006electrical,dahmus2004environmental}. Eliminating it
requires no capital equipment and no process change. It requires only
knowing when a machine is idle, and for how long it will remain so.

Neither is known in an existing facility. Three difficulties compound.

\textbf{There are no labels.} A per-timestamp record of which state each
machine occupied would reduce state identification to a standard
classification problem, but obtaining one means instrumenting every
machine's contactor and control system, or having an operator observe
each machine for months. Neither is affordable at facility scale and
neither transfers to a new machine without repeating the effort. Public
industrial benchmarks reflect this: IMDELD
\cite{martins2018imdeld,imdeld2018dataset}, the benchmark used
throughout this paper, publishes $1$\,Hz electrical measurements from
eight heavy machines and no state annotation at all. Any method that
must be validated against ground truth is therefore inapplicable, and
so is any evaluation metric that assumes it.

\textbf{The duration is unknown when the decision must be made.} The
shutdown trade-off is well understood: de-energising saves standby power
but incurs a restart cost, so it pays only when the idle interval
exceeds a break-even duration \cite{mouzon2007operational}. The
operations-research literature solves the resulting scheduling problem
in considerable depth --- and universally assumes the idle interval is
given, either by a deterministic production schedule or by an
analyst-calibrated arrival process. In a retrofit installation the only
available signal is the machine's own measured electrical behaviour, and
nothing in that literature converts one into the other.

\textbf{The two capabilities have matured separately and remain
unconnected.} Non-intrusive load monitoring can infer what a machine is
doing from its electrical signature \cite{hart1992nonintrusive}, but its
objective is retrospective attribution: it does not forecast and it does
not act. Energy-aware scheduling can act optimally given a duration it
cannot obtain. \emph{The former produces labels no controller consumes;
the latter consumes durations no inference procedure supplies.}
Section~\ref{sec:related} develops this gap across four bodies of
literature and states it precisely.

\medskip
\noindent
To close it we propose a \emph{unified data-driven framework for
industrial standby energy elimination}. Rather than treating state
inference, forecasting and control as separate problems, the framework
treats them as one pipeline whose stages are designed against each
other's requirements: the inference stage is constrained to produce a
temporally coherent sequence because duration forecasting is otherwise
ill-posed, and the forecasting stage is constrained to predict a
duration rather than a class because the decision stage consumes a
duration. The framework has three components.

\begin{enumerate}
  \item \textbf{Physics-gated unsupervised state inference.} A Gaussian
  hidden Markov model \cite{rabiner1989tutorial} estimates the emission
  distributions and transition dynamics of the machine's operating
  states directly from electrical measurements, requiring no
  annotation, and an explicit minimum-dwell constraint enforces
  physically admissible state durations at decode time. Because no
  ground-truth labels exist against which the partition could be
  scored, the fit is then \emph{accepted or rejected} against checks
  drawn from sources external to the model: induction-motor
  power-factor and no-load current physics, variance and power
  ordering, agreement with independently derived threshold classifiers,
  and the plant's published shift calendar.

  \item \textbf{Non-productive duration forecasting.} The recent
  inferred-state history is mapped to the duration of the upcoming
  non-productive interval. Predicting duration rather than the next
  state label is what makes the output actionable, since the break-even
  criterion is a statement about how long a machine will remain idle.
  Six trained forecasters --- recurrent, convolutional, transformer and
  tree-based --- are compared under a common protocol \emph{and against
  an untrained persistence reference}, which is the comparison that
  determines whether the stage earns its place at all.

  \item \textbf{Optimisation-based shutdown decision support.} The
  decision is posed as minimisation of total cost --- standby energy,
  restart, and production delay --- subject to thermal cool-down and
  restart-frequency constraints, in the spirit of
  \cite{mouzon2007operational} but with the idle duration supplied by
  component~(ii) instead of by a known schedule. The break-even
  duration is therefore a \emph{derived} quantity of the optimisation
  rather than an operator-supplied threshold, and because the duration
  is unknown at decision time the policy is benchmarked against the
  $2$-competitive forecast-free rent-or-buy rule
  \cite{karlin1994competitive} rather than against trivial policies.
\end{enumerate}

\medskip
\noindent
\textbf{Scope of the claim.} We claim no novelty in the Gaussian hidden
Markov model, in any of the forecasting architectures, or in the
break-even criterion, each of which is established. The contribution is
the framework that connects them --- specifically, the demonstration
that label-free state inference can be made trustworthy enough, and
temporally coherent enough, to serve as the input to an economic control
decision on real industrial data --- together with an evaluation
honest enough to identify which of its own components do not earn their
place.

\subsection*{Contributions}

\begin{itemize}
  \item \textbf{A physics accept/reject loop, and the measurement that
  shows it is load-bearing.} Fitting the labeller at ten seeds on
  identical preprocessed data moves the headline standby total by a
  coefficient of variation of $33.7\%$, and in three seeds of ten the
  cluster named \textsc{standby} is physically a different state
  --- once a loaded running state at $13$\,kW, twice the machine
  switched off at under $1$\,W. The physics battery scores exactly those
  three below threshold, and conditioning on a passing score collapses
  the spread to $3.0\%$. \emph{The pipeline is reproducible because of
  the physics validation, not despite the clustering}
  (Sections~\ref{sec:method:accept_reject},
  \ref{sec:significance:seeds}).

  \item \textbf{An account of what temporal coherence actually
  requires.} Viterbi decoding over a learned transition matrix reduces
  flicker substantially but --- contrary to the usual assumption ---
  cannot eliminate it, and neither can strengthening the transition
  prior. The penalty a transition matrix can impose on a state change is
  bounded by $\log(p_{\text{self}}/p_{\text{switch}})$, measured here at
  $9.23$ nats, whereas full-covariance Gaussian emissions fitted to
  low-noise electrical data differ by a median of $27.8$ nats and a
  $90$th percentile of $41\,701$ nats between adjacent states. At
  $55.7\%$ of timesteps the emission term is arithmetically unable to be
  overruled. We show the remedy is an explicit minimum-dwell constraint,
  that it removes flicker entirely at a cost of $0.19\%$ of rows, and
  that it does so on all eight machines
  (Sections~\ref{sec:flicker}, \ref{sec:multimachine:general}).

  \item \textbf{A validation methodology for state partitions in the
  absence of ground truth.} Circuit-level physics, model-level
  diagnostics, and independently derived Otsu-thresholded classifiers
  \cite{otsu1979threshold} are combined by consensus
  \cite{strehl2002cluster}, none of them using the model's own internal
  statistics. Applied across eight machines the battery accepts the four
  motor-driven ones and rejects the four auxiliary ones, on which there
  is no energised-idle state to find --- which is the behaviour that
  makes it a validation procedure rather than a formality
  (Sections~\ref{sec:validation}, \ref{sec:multimachine:states}).

  \item \textbf{A forecasting comparison anchored on not forecasting,
  and a negative result reported as such.} Under one protocol with
  seeds, an untrained persistence reference beats all five neural
  sequence models, the originally proposed sequence-to-sequence LSTM by
  the largest margin of the seven ($7.95$\,s of per-window MAE,
  $p=5.7\times10^{-44}$). A gradient-boosted tree is the only trained
  forecaster that clears the reference, and it is the forecasting
  component of the framework as reported here
  (Sections~\ref{sec:forecasting}, \ref{sec:significance:forecast}).

  \item \textbf{A decision layer in which the break-even duration is
  derived, with an invariance result and a competitive benchmark.} We
  replace the fixed-threshold rule with a constrained cost
  minimisation, show that the reduced problem admits an exact offline
  optimum, and measure the restart coefficients from the record instead
  of assuming them --- finding the studied machine has no measurable
  electrical restart penalty, so its restart cost is a production-delay
  cost. We show analytically and numerically that the break-even
  duration is \emph{independent of electricity tariff} whenever the
  restart cost is purely electrical, and report the decision surface
  over the axes that do carry information. The policy beats the fixed
  threshold it replaces by $+68.9$\,USD/yr ($p=0.001$); it is not shown
  to beat the plant's own behaviour, and we say so
  (Sections~\ref{sec:decision}, \ref{sec:significance:decision}).

  \item \textbf{End-to-end evaluation with ablation, significance
  testing and transfer.} Every stage is evaluated across eight machines
  and two sites, each component's marginal contribution is priced by
  ablation, every headline difference carries a dependence-aware
  interval and a pre-registered practical threshold, and the two designs
  that cannot detect anything are identified as such rather than
  reported as null results (Sections~\ref{sec:multimachine}--\ref{sec:ablation}).

  \item \textbf{Reproducibility.} The full pipeline --- preprocessing,
  inference, validation harness, forecasting, decision layer and every
  experiment reported here --- is released as open source, and all
  results are obtained from public benchmark data.
\end{itemize}

\subsection*{Organisation}

Section~\ref{sec:related} reviews related work across machine state
identification, non-intrusive load monitoring, industrial energy
optimisation and predictive shutdown, and states the resulting gap.
Section~\ref{sec:method} presents the framework.
Section~\ref{sec:states} reports the state layer, including the
comparison against clustering baselines, the flicker analysis, and the
ground-truth-free validation. Section~\ref{sec:forecasting} reports the
forecasting comparison and Section~\ref{sec:significance} the
statistical analysis on which every subsequent claim depends.
Section~\ref{sec:decision} develops the decision layer.
Sections~\ref{sec:multimachine} and~\ref{sec:crossmachine} extend the
evaluation to all eight machines, to transfer between them, and to a
second site. Section~\ref{sec:explain} attributes what each model uses,
Section~\ref{sec:ablation} prices each component, and
Section~\ref{sec:limitations} states the boundaries of the evidence
before Section~\ref{sec:conclusion} concludes.

%  Bibliography: paper/references.bib
% =====================================================================

\section{Introduction}
\label{sec:intro}

Electric motor systems account for the largest single share of
industrial electricity use, and a substantial part of that share is
drawn while the machine is energised but producing nothing. An induction
motor left running between production runs continues to draw magnetising
current: real power falls to a fraction of its loaded value, but it does
not fall to zero, and the machine's power factor collapses. ISO~14955-1
\cite{iso14955} names this operating mode --- ready for production,
performing no value-adding function --- and studies of discrete-part
manufacturing find that such load-independent consumption frequently
exceeds the energy actually delivered to the workpiece
\cite{gutowski2006electrical,dahmus2004environmental}. Eliminating it
requires no capital equipment and no process change. It requires only
knowing when a machine is idle, and for how long it will remain so.

Neither is known in an existing facility. Three difficulties compound.

\textbf{There are no labels.} A per-timestamp record of which state each
machine occupied would reduce state identification to a standard
classification problem, but obtaining one means instrumenting every
machine's contactor and control system, or having an operator observe
each machine for months. Neither is affordable at facility scale and
neither transfers to a new machine without repeating the effort. Public
industrial benchmarks reflect this: IMDELD
\cite{martins2018imdeld,imdeld2018dataset}, the benchmark used
throughout this paper, publishes $1$\,Hz electrical measurements from
eight heavy machines and no state annotation at all. Any method that
must be validated against ground truth is therefore inapplicable, and
so is any evaluation metric that assumes it.

\textbf{The duration is unknown when the decision must be made.} The
shutdown trade-off is well understood: de-energising saves standby power
but incurs a restart cost, so it pays only when the idle interval
exceeds a break-even duration \cite{mouzon2007operational}. The
operations-research literature solves the resulting scheduling problem
in considerable depth --- and universally assumes the idle interval is
given, either by a deterministic production schedule or by an
analyst-calibrated arrival process. In a retrofit installation the only
available signal is the machine's own measured electrical behaviour, and
nothing in that literature converts one into the other.

\textbf{The two capabilities have matured separately and remain
unconnected.} Non-intrusive load monitoring can infer what a machine is
doing from its electrical signature \cite{hart1992nonintrusive}, but its
objective is retrospective attribution: it does not forecast and it does
not act. Energy-aware scheduling can act optimally given a duration it
cannot obtain. \emph{The former produces labels no controller consumes;
the latter consumes durations no inference procedure supplies.}
Section~\ref{sec:related} develops this gap across four bodies of
literature and states it precisely.

\medskip
\noindent
To close it we propose a \emph{unified data-driven framework for
industrial standby energy elimination}. Rather than treating state
inference, forecasting and control as separate problems, the framework
treats them as one pipeline whose stages are designed against each
other's requirements: the inference stage is constrained to produce a
temporally coherent sequence because duration forecasting is otherwise
ill-posed, and the forecasting stage is constrained to predict a
duration rather than a class because the decision stage consumes a
duration. The framework has three components.

\begin{enumerate}
  \item \textbf{Physics-gated unsupervised state inference.} A Gaussian
  hidden Markov model \cite{rabiner1989tutorial} estimates the emission
  distributions and transition dynamics of the machine's operating
  states directly from electrical measurements, requiring no
  annotation, and an explicit minimum-dwell constraint enforces
  physically admissible state durations at decode time. Because no
  ground-truth labels exist against which the partition could be
  scored, the fit is then \emph{accepted or rejected} against checks
  drawn from sources external to the model: induction-motor
  power-factor and no-load current physics, variance and power
  ordering, agreement with independently derived threshold classifiers,
  and the plant's published shift calendar.

  \item \textbf{Non-productive duration forecasting.} The recent
  inferred-state history is mapped to the duration of the upcoming
  non-productive interval. Predicting duration rather than the next
  state label is what makes the output actionable, since the break-even
  criterion is a statement about how long a machine will remain idle.
  Six trained forecasters --- recurrent, convolutional, transformer and
  tree-based --- are compared under a common protocol \emph{and against
  an untrained persistence reference}, which is the comparison that
  determines whether the stage earns its place at all.

  \item \textbf{Optimisation-based shutdown decision support.} The
  decision is posed as minimisation of total cost --- standby energy,
  restart, and production delay --- subject to thermal cool-down and
  restart-frequency constraints, in the spirit of
  \cite{mouzon2007operational} but with the idle duration supplied by
  component~(ii) instead of by a known schedule. The break-even
  duration is therefore a \emph{derived} quantity of the optimisation
  rather than an operator-supplied threshold, and because the duration
  is unknown at decision time the policy is benchmarked against the
  $2$-competitive forecast-free rent-or-buy rule
  \cite{karlin1994competitive} rather than against trivial policies.
\end{enumerate}

\medskip
\noindent
\textbf{Scope of the claim.} We claim no novelty in the Gaussian hidden
Markov model, in any of the forecasting architectures, or in the
break-even criterion, each of which is established. The contribution is
the framework that connects them --- specifically, the demonstration
that label-free state inference can be made trustworthy enough, and
temporally coherent enough, to serve as the input to an economic control
decision on real industrial data --- together with an evaluation
honest enough to identify which of its own components do not earn their
place.

\subsection*{Contributions}

\begin{itemize}
  \item \textbf{A physics accept/reject loop, and the measurement that
  shows it is load-bearing.} Fitting the labeller at ten seeds on
  identical preprocessed data moves the headline standby total by a
  coefficient of variation of $33.7\%$, and in three seeds of ten the
  cluster named \textsc{standby} is physically a different state
  --- once a loaded running state at $13$\,kW, twice the machine
  switched off at under $1$\,W. The physics battery scores exactly those
  three below threshold, and conditioning on a passing score collapses
  the spread to $3.0\%$. \emph{The pipeline is reproducible because of
  the physics validation, not despite the clustering}
  (Sections~\ref{sec:method:accept_reject},
  \ref{sec:significance:seeds}).

  \item \textbf{An account of what temporal coherence actually
  requires.} Viterbi decoding over a learned transition matrix reduces
  flicker substantially but --- contrary to the usual assumption ---
  cannot eliminate it, and neither can strengthening the transition
  prior. The penalty a transition matrix can impose on a state change is
  bounded by $\log(p_{\text{self}}/p_{\text{switch}})$, measured here at
  $9.23$ nats, whereas full-covariance Gaussian emissions fitted to
  low-noise electrical data differ by a median of $27.8$ nats and a
  $90$th percentile of $41\,701$ nats between adjacent states. At
  $55.7\%$ of timesteps the emission term is arithmetically unable to be
  overruled. We show the remedy is an explicit minimum-dwell constraint,
  that it removes flicker entirely at a cost of $0.19\%$ of rows, and
  that it does so on all eight machines
  (Sections~\ref{sec:flicker}, \ref{sec:multimachine:general}).

  \item \textbf{A validation methodology for state partitions in the
  absence of ground truth.} Circuit-level physics, model-level
  diagnostics, and independently derived Otsu-thresholded classifiers
  \cite{otsu1979threshold} are combined by consensus
  \cite{strehl2002cluster}, none of them using the model's own internal
  statistics. Applied across eight machines the battery accepts the four
  motor-driven ones and rejects the four auxiliary ones, on which there
  is no energised-idle state to find --- which is the behaviour that
  makes it a validation procedure rather than a formality
  (Sections~\ref{sec:validation}, \ref{sec:multimachine:states}).

  \item \textbf{A forecasting comparison anchored on not forecasting,
  and a negative result reported as such.} Under one protocol with
  seeds, an untrained persistence reference beats all five neural
  sequence models, the originally proposed sequence-to-sequence LSTM by
  the largest margin of the seven ($7.95$\,s of per-window MAE,
  $p=5.7\times10^{-44}$). A gradient-boosted tree is the only trained
  forecaster that clears the reference, and it is the forecasting
  component of the framework as reported here
  (Sections~\ref{sec:forecasting}, \ref{sec:significance:forecast}).

  \item \textbf{A decision layer in which the break-even duration is
  derived, with an invariance result and a competitive benchmark.} We
  replace the fixed-threshold rule with a constrained cost
  minimisation, show that the reduced problem admits an exact offline
  optimum, and measure the restart coefficients from the record instead
  of assuming them --- finding the studied machine has no measurable
  electrical restart penalty, so its restart cost is a production-delay
  cost. We show analytically and numerically that the break-even
  duration is \emph{independent of electricity tariff} whenever the
  restart cost is purely electrical, and report the decision surface
  over the axes that do carry information. The policy beats the fixed
  threshold it replaces by $+68.9$\,USD/yr ($p=0.001$); it is not shown
  to beat the plant's own behaviour, and we say so
  (Sections~\ref{sec:decision}, \ref{sec:significance:decision}).

  \item \textbf{End-to-end evaluation with ablation, significance
  testing and transfer.} Every stage is evaluated across eight machines
  and two sites, each component's marginal contribution is priced by
  ablation, every headline difference carries a dependence-aware
  interval and a pre-registered practical threshold, and the two designs
  that cannot detect anything are identified as such rather than
  reported as null results (Sections~\ref{sec:multimachine}--\ref{sec:ablation}).

  \item \textbf{Reproducibility.} The full pipeline --- preprocessing,
  inference, validation harness, forecasting, decision layer and every
  experiment reported here --- is released as open source, and all
  results are obtained from public benchmark data.
\end{itemize}

\subsection*{Organisation}

Section~\ref{sec:related} reviews related work across machine state
identification, non-intrusive load monitoring, industrial energy
optimisation and predictive shutdown, and states the resulting gap.
Section~\ref{sec:method} presents the framework.
Section~\ref{sec:states} reports the state layer, including the
comparison against clustering baselines, the flicker analysis, and the
ground-truth-free validation. Section~\ref{sec:forecasting} reports the
forecasting comparison and Section~\ref{sec:significance} the
statistical analysis on which every subsequent claim depends.
Section~\ref{sec:decision} develops the decision layer.
Sections~\ref{sec:multimachine} and~\ref{sec:crossmachine} extend the
evaluation to all eight machines, to transfer between them, and to a
second site. Section~\ref{sec:explain} attributes what each model uses,
Section~\ref{sec:ablation} prices each component, and
Section~\ref{sec:limitations} states the boundaries of the evidence
before Section~\ref{sec:conclusion} concludes.

%  Bibliography: paper/references.bib
% =====================================================================

\section{Introduction}
\label{sec:intro}

Electric motor systems account for the largest single share of
industrial electricity use, and a substantial part of that share is
drawn while the machine is energised but producing nothing. An induction
motor left running between production runs continues to draw magnetising
current: real power falls to a fraction of its loaded value, but it does
not fall to zero, and the machine's power factor collapses. ISO~14955-1
\cite{iso14955} names this operating mode --- ready for production,
performing no value-adding function --- and studies of discrete-part
manufacturing find that such load-independent consumption frequently
exceeds the energy actually delivered to the workpiece
\cite{gutowski2006electrical,dahmus2004environmental}. Eliminating it
requires no capital equipment and no process change. It requires only
knowing when a machine is idle, and for how long it will remain so.

Neither is known in an existing facility. Three difficulties compound.

\textbf{There are no labels.} A per-timestamp record of which state each
machine occupied would reduce state identification to a standard
classification problem, but obtaining one means instrumenting every
machine's contactor and control system, or having an operator observe
each machine for months. Neither is affordable at facility scale and
neither transfers to a new machine without repeating the effort. Public
industrial benchmarks reflect this: IMDELD
\cite{martins2018imdeld,imdeld2018dataset}, the benchmark used
throughout this paper, publishes $1$\,Hz electrical measurements from
eight heavy machines and no state annotation at all. Any method that
must be validated against ground truth is therefore inapplicable, and
so is any evaluation metric that assumes it.

\textbf{The duration is unknown when the decision must be made.} The
shutdown trade-off is well understood: de-energising saves standby power
but incurs a restart cost, so it pays only when the idle interval
exceeds a break-even duration \cite{mouzon2007operational}. The
operations-research literature solves the resulting scheduling problem
in considerable depth --- and universally assumes the idle interval is
given, either by a deterministic production schedule or by an
analyst-calibrated arrival process. In a retrofit installation the only
available signal is the machine's own measured electrical behaviour, and
nothing in that literature converts one into the other.

\textbf{The two capabilities have matured separately and remain
unconnected.} Non-intrusive load monitoring can infer what a machine is
doing from its electrical signature \cite{hart1992nonintrusive}, but its
objective is retrospective attribution: it does not forecast and it does
not act. Energy-aware scheduling can act optimally given a duration it
cannot obtain. \emph{The former produces labels no controller consumes;
the latter consumes durations no inference procedure supplies.}
Section~\ref{sec:related} develops this gap across four bodies of
literature and states it precisely.

\medskip
\noindent
To close it we propose a \emph{unified data-driven framework for
industrial standby energy elimination}. Rather than treating state
inference, forecasting and control as separate problems, the framework
treats them as one pipeline whose stages are designed against each
other's requirements: the inference stage is constrained to produce a
temporally coherent sequence because duration forecasting is otherwise
ill-posed, and the forecasting stage is constrained to predict a
duration rather than a class because the decision stage consumes a
duration. The framework has three components.

\begin{enumerate}
  \item \textbf{Physics-gated unsupervised state inference.} A Gaussian
  hidden Markov model \cite{rabiner1989tutorial} estimates the emission
  distributions and transition dynamics of the machine's operating
  states directly from electrical measurements, requiring no
  annotation, and an explicit minimum-dwell constraint enforces
  physically admissible state durations at decode time. Because no
  ground-truth labels exist against which the partition could be
  scored, the fit is then \emph{accepted or rejected} against checks
  drawn from sources external to the model: induction-motor
  power-factor and no-load current physics, variance and power
  ordering, agreement with independently derived threshold classifiers,
  and the plant's published shift calendar.

  \item \textbf{Non-productive duration forecasting.} The recent
  inferred-state history is mapped to the duration of the upcoming
  non-productive interval. Predicting duration rather than the next
  state label is what makes the output actionable, since the break-even
  criterion is a statement about how long a machine will remain idle.
  Six trained forecasters --- recurrent, convolutional, transformer and
  tree-based --- are compared under a common protocol \emph{and against
  an untrained persistence reference}, which is the comparison that
  determines whether the stage earns its place at all.

  \item \textbf{Optimisation-based shutdown decision support.} The
  decision is posed as minimisation of total cost --- standby energy,
  restart, and production delay --- subject to thermal cool-down and
  restart-frequency constraints, in the spirit of
  \cite{mouzon2007operational} but with the idle duration supplied by
  component~(ii) instead of by a known schedule. The break-even
  duration is therefore a \emph{derived} quantity of the optimisation
  rather than an operator-supplied threshold, and because the duration
  is unknown at decision time the policy is benchmarked against the
  $2$-competitive forecast-free rent-or-buy rule
  \cite{karlin1994competitive} rather than against trivial policies.
\end{enumerate}

\medskip
\noindent
\textbf{Scope of the claim.} We claim no novelty in the Gaussian hidden
Markov model, in any of the forecasting architectures, or in the
break-even criterion, each of which is established. The contribution is
the framework that connects them --- specifically, the demonstration
that label-free state inference can be made trustworthy enough, and
temporally coherent enough, to serve as the input to an economic control
decision on real industrial data --- together with an evaluation
honest enough to identify which of its own components do not earn their
place.

\subsection*{Contributions}

\begin{itemize}
  \item \textbf{A physics accept/reject loop, and the measurement that
  shows it is load-bearing.} Fitting the labeller at ten seeds on
  identical preprocessed data moves the headline standby total by a
  coefficient of variation of $33.7\%$, and in three seeds of ten the
  cluster named \textsc{standby} is physically a different state
  --- once a loaded running state at $13$\,kW, twice the machine
  switched off at under $1$\,W. The physics battery scores exactly those
  three below threshold, and conditioning on a passing score collapses
  the spread to $3.0\%$. \emph{The pipeline is reproducible because of
  the physics validation, not despite the clustering}
  (Sections~\ref{sec:method:accept_reject},
  \ref{sec:significance:seeds}).

  \item \textbf{An account of what temporal coherence actually
  requires.} Viterbi decoding over a learned transition matrix reduces
  flicker substantially but --- contrary to the usual assumption ---
  cannot eliminate it, and neither can strengthening the transition
  prior. The penalty a transition matrix can impose on a state change is
  bounded by $\log(p_{\text{self}}/p_{\text{switch}})$, measured here at
  $9.23$ nats, whereas full-covariance Gaussian emissions fitted to
  low-noise electrical data differ by a median of $27.8$ nats and a
  $90$th percentile of $41\,701$ nats between adjacent states. At
  $55.7\%$ of timesteps the emission term is arithmetically unable to be
  overruled. We show the remedy is an explicit minimum-dwell constraint,
  that it removes flicker entirely at a cost of $0.19\%$ of rows, and
  that it does so on all eight machines
  (Sections~\ref{sec:flicker}, \ref{sec:multimachine:general}).

  \item \textbf{A validation methodology for state partitions in the
  absence of ground truth.} Circuit-level physics, model-level
  diagnostics, and independently derived Otsu-thresholded classifiers
  \cite{otsu1979threshold} are combined by consensus
  \cite{strehl2002cluster}, none of them using the model's own internal
  statistics. Applied across eight machines the battery accepts the four
  motor-driven ones and rejects the four auxiliary ones, on which there
  is no energised-idle state to find --- which is the behaviour that
  makes it a validation procedure rather than a formality
  (Sections~\ref{sec:validation}, \ref{sec:multimachine:states}).

  \item \textbf{A forecasting comparison anchored on not forecasting,
  and a negative result reported as such.} Under one protocol with
  seeds, an untrained persistence reference beats all five neural
  sequence models, the originally proposed sequence-to-sequence LSTM by
  the largest margin of the seven ($7.95$\,s of per-window MAE,
  $p=5.7\times10^{-44}$). A gradient-boosted tree is the only trained
  forecaster that clears the reference, and it is the forecasting
  component of the framework as reported here
  (Sections~\ref{sec:forecasting}, \ref{sec:significance:forecast}).

  \item \textbf{A decision layer in which the break-even duration is
  derived, with an invariance result and a competitive benchmark.} We
  replace the fixed-threshold rule with a constrained cost
  minimisation, show that the reduced problem admits an exact offline
  optimum, and measure the restart coefficients from the record instead
  of assuming them --- finding the studied machine has no measurable
  electrical restart penalty, so its restart cost is a production-delay
  cost. We show analytically and numerically that the break-even
  duration is \emph{independent of electricity tariff} whenever the
  restart cost is purely electrical, and report the decision surface
  over the axes that do carry information. The policy beats the fixed
  threshold it replaces by $+68.9$\,USD/yr ($p=0.001$); it is not shown
  to beat the plant's own behaviour, and we say so
  (Sections~\ref{sec:decision}, \ref{sec:significance:decision}).

  \item \textbf{End-to-end evaluation with ablation, significance
  testing and transfer.} Every stage is evaluated across eight machines
  and two sites, each component's marginal contribution is priced by
  ablation, every headline difference carries a dependence-aware
  interval and a pre-registered practical threshold, and the two designs
  that cannot detect anything are identified as such rather than
  reported as null results (Sections~\ref{sec:multimachine}--\ref{sec:ablation}).

  \item \textbf{Reproducibility.} The full pipeline --- preprocessing,
  inference, validation harness, forecasting, decision layer and every
  experiment reported here --- is released as open source, and all
  results are obtained from public benchmark data.
\end{itemize}

\subsection*{Organisation}

Section~\ref{sec:related} reviews related work across machine state
identification, non-intrusive load monitoring, industrial energy
optimisation and predictive shutdown, and states the resulting gap.
Section~\ref{sec:method} presents the framework.
Section~\ref{sec:states} reports the state layer, including the
comparison against clustering baselines, the flicker analysis, and the
ground-truth-free validation. Section~\ref{sec:forecasting} reports the
forecasting comparison and Section~\ref{sec:significance} the
statistical analysis on which every subsequent claim depends.
Section~\ref{sec:decision} develops the decision layer.
Sections~\ref{sec:multimachine} and~\ref{sec:crossmachine} extend the
evaluation to all eight machines, to transfer between them, and to a
second site. Section~\ref{sec:explain} attributes what each model uses,
Section~\ref{sec:ablation} prices each component, and
Section~\ref{sec:limitations} states the boundaries of the evidence
before Section~\ref{sec:conclusion} concludes.

%  Bibliography: paper/references.bib
% =====================================================================

\section{Introduction}
\label{sec:intro}

Electric motor systems account for the largest single share of
industrial electricity use, and a substantial part of that share is
drawn while the machine is energised but producing nothing. An induction
motor left running between production runs continues to draw magnetising
current: real power falls to a fraction of its loaded value, but it does
not fall to zero, and the machine's power factor collapses. ISO~14955-1
\cite{iso14955} names this operating mode --- ready for production,
performing no value-adding function --- and studies of discrete-part
manufacturing find that such load-independent consumption frequently
exceeds the energy actually delivered to the workpiece
\cite{gutowski2006electrical,dahmus2004environmental}. Eliminating it
requires no capital equipment and no process change. It requires only
knowing when a machine is idle, and for how long it will remain so.

Neither is known in an existing facility. Three difficulties compound.

\textbf{There are no labels.} A per-timestamp record of which state each
machine occupied would reduce state identification to a standard
classification problem, but obtaining one means instrumenting every
machine's contactor and control system, or having an operator observe
each machine for months. Neither is affordable at facility scale and
neither transfers to a new machine without repeating the effort. Public
industrial benchmarks reflect this: IMDELD
\cite{martins2018imdeld,imdeld2018dataset}, the benchmark used
throughout this paper, publishes $1$\,Hz electrical measurements from
eight heavy machines and no state annotation at all. Any method that
must be validated against ground truth is therefore inapplicable, and
so is any evaluation metric that assumes it.

\textbf{The duration is unknown when the decision must be made.} The
shutdown trade-off is well understood: de-energising saves standby power
but incurs a restart cost, so it pays only when the idle interval
exceeds a break-even duration \cite{mouzon2007operational}. The
operations-research literature solves the resulting scheduling problem
in considerable depth --- and universally assumes the idle interval is
given, either by a deterministic production schedule or by an
analyst-calibrated arrival process. In a retrofit installation the only
available signal is the machine's own measured electrical behaviour, and
nothing in that literature converts one into the other.

\textbf{The two capabilities have matured separately and remain
unconnected.} Non-intrusive load monitoring can infer what a machine is
doing from its electrical signature \cite{hart1992nonintrusive}, but its
objective is retrospective attribution: it does not forecast and it does
not act. Energy-aware scheduling can act optimally given a duration it
cannot obtain. \emph{The former produces labels no controller consumes;
the latter consumes durations no inference procedure supplies.}
Section~\ref{sec:related} develops this gap across four bodies of
literature and states it precisely.

\medskip
\noindent
To close it we propose a \emph{unified data-driven framework for
industrial standby energy elimination}. Rather than treating state
inference, forecasting and control as separate problems, the framework
treats them as one pipeline whose stages are designed against each
other's requirements: the inference stage is constrained to produce a
temporally coherent sequence because duration forecasting is otherwise
ill-posed, and the forecasting stage is constrained to predict a
duration rather than a class because the decision stage consumes a
duration. The framework has three components.

\begin{enumerate}
  \item \textbf{Physics-gated unsupervised state inference.} A Gaussian
  hidden Markov model \cite{rabiner1989tutorial} estimates the emission
  distributions and transition dynamics of the machine's operating
  states directly from electrical measurements, requiring no
  annotation, and an explicit minimum-dwell constraint enforces
  physically admissible state durations at decode time. Because no
  ground-truth labels exist against which the partition could be
  scored, the fit is then \emph{accepted or rejected} against checks
  drawn from sources external to the model: induction-motor
  power-factor and no-load current physics, variance and power
  ordering, agreement with independently derived threshold classifiers,
  and the plant's published shift calendar.

  \item \textbf{Non-productive duration forecasting.} The recent
  inferred-state history is mapped to the duration of the upcoming
  non-productive interval. Predicting duration rather than the next
  state label is what makes the output actionable, since the break-even
  criterion is a statement about how long a machine will remain idle.
  Six trained forecasters --- recurrent, convolutional, transformer and
  tree-based --- are compared under a common protocol \emph{and against
  an untrained persistence reference}, which is the comparison that
  determines whether the stage earns its place at all.

  \item \textbf{Optimisation-based shutdown decision support.} The
  decision is posed as minimisation of total cost --- standby energy,
  restart, and production delay --- subject to thermal cool-down and
  restart-frequency constraints, in the spirit of
  \cite{mouzon2007operational} but with the idle duration supplied by
  component~(ii) instead of by a known schedule. The break-even
  duration is therefore a \emph{derived} quantity of the optimisation
  rather than an operator-supplied threshold, and because the duration
  is unknown at decision time the policy is benchmarked against the
  $2$-competitive forecast-free rent-or-buy rule
  \cite{karlin1994competitive} rather than against trivial policies.
\end{enumerate}

\medskip
\noindent
\textbf{Scope of the claim.} We claim no novelty in the Gaussian hidden
Markov model, in any of the forecasting architectures, or in the
break-even criterion, each of which is established. The contribution is
the framework that connects them --- specifically, the demonstration
that label-free state inference can be made trustworthy enough, and
temporally coherent enough, to serve as the input to an economic control
decision on real industrial data --- together with an evaluation
honest enough to identify which of its own components do not earn their
place.

\subsection*{Contributions}

\begin{itemize}
  \item \textbf{A physics accept/reject loop, and the measurement that
  shows it is load-bearing.} Fitting the labeller at ten seeds on
  identical preprocessed data moves the headline standby total by a
  coefficient of variation of $33.7\%$, and in three seeds of ten the
  cluster named \textsc{standby} is physically a different state
  --- once a loaded running state at $13$\,kW, twice the machine
  switched off at under $1$\,W. The physics battery scores exactly those
  three below threshold, and conditioning on a passing score collapses
  the spread to $3.0\%$. \emph{The pipeline is reproducible because of
  the physics validation, not despite the clustering}
  (Sections~\ref{sec:method:accept_reject},
  \ref{sec:significance:seeds}).

  \item \textbf{An account of what temporal coherence actually
  requires.} Viterbi decoding over a learned transition matrix reduces
  flicker substantially but --- contrary to the usual assumption ---
  cannot eliminate it, and neither can strengthening the transition
  prior. The penalty a transition matrix can impose on a state change is
  bounded by $\log(p_{\text{self}}/p_{\text{switch}})$, measured here at
  $9.23$ nats, whereas full-covariance Gaussian emissions fitted to
  low-noise electrical data differ by a median of $27.8$ nats and a
  $90$th percentile of $41\,701$ nats between adjacent states. At
  $55.7\%$ of timesteps the emission term is arithmetically unable to be
  overruled. We show the remedy is an explicit minimum-dwell constraint,
  that it removes flicker entirely at a cost of $0.19\%$ of rows, and
  that it does so on all eight machines
  (Sections~\ref{sec:flicker}, \ref{sec:multimachine:general}).

  \item \textbf{A validation methodology for state partitions in the
  absence of ground truth.} Circuit-level physics, model-level
  diagnostics, and independently derived Otsu-thresholded classifiers
  \cite{otsu1979threshold} are combined by consensus
  \cite{strehl2002cluster}, none of them using the model's own internal
  statistics. Applied across eight machines the battery accepts the four
  motor-driven ones and rejects the four auxiliary ones, on which there
  is no energised-idle state to find --- which is the behaviour that
  makes it a validation procedure rather than a formality
  (Sections~\ref{sec:validation}, \ref{sec:multimachine:states}).

  \item \textbf{A forecasting comparison anchored on not forecasting,
  and a negative result reported as such.} Under one protocol with
  seeds, an untrained persistence reference beats all five neural
  sequence models, the originally proposed sequence-to-sequence LSTM by
  the largest margin of the seven ($7.95$\,s of per-window MAE,
  $p=5.7\times10^{-44}$). A gradient-boosted tree is the only trained
  forecaster that clears the reference, and it is the forecasting
  component of the framework as reported here
  (Sections~\ref{sec:forecasting}, \ref{sec:significance:forecast}).

  \item \textbf{A decision layer in which the break-even duration is
  derived, with an invariance result and a competitive benchmark.} We
  replace the fixed-threshold rule with a constrained cost
  minimisation, show that the reduced problem admits an exact offline
  optimum, and measure the restart coefficients from the record instead
  of assuming them --- finding the studied machine has no measurable
  electrical restart penalty, so its restart cost is a production-delay
  cost. We show analytically and numerically that the break-even
  duration is \emph{independent of electricity tariff} whenever the
  restart cost is purely electrical, and report the decision surface
  over the axes that do carry information. The policy beats the fixed
  threshold it replaces by $+68.9$\,USD/yr ($p=0.001$); it is not shown
  to beat the plant's own behaviour, and we say so
  (Sections~\ref{sec:decision}, \ref{sec:significance:decision}).

  \item \textbf{End-to-end evaluation with ablation, significance
  testing and transfer.} Every stage is evaluated across eight machines
  and two sites, each component's marginal contribution is priced by
  ablation, every headline difference carries a dependence-aware
  interval and a pre-registered practical threshold, and the two designs
  that cannot detect anything are identified as such rather than
  reported as null results (Sections~\ref{sec:multimachine}--\ref{sec:ablation}).

  \item \textbf{Reproducibility.} The full pipeline --- preprocessing,
  inference, validation harness, forecasting, decision layer and every
  experiment reported here --- is released as open source, and all
  results are obtained from public benchmark data.
\end{itemize}

\subsection*{Organisation}

Section~\ref{sec:related} reviews related work across machine state
identification, non-intrusive load monitoring, industrial energy
optimisation and predictive shutdown, and states the resulting gap.
Section~\ref{sec:method} presents the framework.
Section~\ref{sec:states} reports the state layer, including the
comparison against clustering baselines, the flicker analysis, and the
ground-truth-free validation. Section~\ref{sec:forecasting} reports the
forecasting comparison and Section~\ref{sec:significance} the
statistical analysis on which every subsequent claim depends.
Section~\ref{sec:decision} develops the decision layer.
Sections~\ref{sec:multimachine} and~\ref{sec:crossmachine} extend the
evaluation to all eight machines, to transfer between them, and to a
second site. Section~\ref{sec:explain} attributes what each model uses,
Section~\ref{sec:ablation} prices each component, and
Section~\ref{sec:limitations} states the boundaries of the evidence
before Section~\ref{sec:conclusion} concludes.
